\documentclass[11pt]{article}
\usepackage{amsmath,amssymb,amsthm,geometry}
\geometry{margin=.9in}\setlength{\emergencystretch}{1.5em}
\newtheorem{theorem}{Theorem}\newtheorem{lemma}{Lemma}
\newcommand{\dd}{\,d}\newcommand{\E}{\mathbb E}
\title{The trained half-mass band returns before weak learning}
\author{AN04 proof increment}\date{Version 1}
\begin{document}\maketitle

\paragraph{Definitions and domains.}
Fix the HBD critical Gaussian family $S_0=cq^{1/\omega}$, aligned
balanced isotropic continuum, $c>0$, $\rho\in[13/20,7/10]$,
$\lambda=\rho^2$, $\omega=(1-\lambda)/2$, and fixed $\zeta>a_\rho$.
Use normalized time, $t_h$ the first right-ancestry half mass, and
$s=t-t_h$. Let $B$ be the fixed positive-length compact interior Q1
band of HBD main, chosen inactive at $t_h$.
The probe is $\xi=(\sin(q\zeta),-\cos(q\zeta))$. In any reached
bounded strong chart $x=\theta/q$ it is active precisely if $x<\zeta$.
Weak learning is measured by P learning,
$\mathsf m_2=\E_{P_2}[D_2 f_2(D)]/(\rho^2+q^2)$,
where $D$ is the data vector. A positive weak radial mass in a
chosen chart is a different event. Write $t_\eta$ for the first
post-half-mass time $\mathsf m_2\ge\eta$ for fixed $\eta>0$, with
infinite hitting time allowed.

\begin{theorem}[Trained flow]\label{an04brwv1:return}
There exist $T_B<\infty$ and $d_B>0$, independent of $q$, such that,
for small $q$, every label in $B$ has $x(t_h+T_B)\le\zeta-d_B$.
More generally this holds uniformly on $T_B\le s\le T$ for each
fixed finite $T\ge T_B$, taking $q$ small after $T$. Moreover there
is $k>0$ such that
\[
 \sup_{0\le s\le k\sqrt{\log(1/q)}}|\mathsf m_2(t_h+s)|\longrightarrow0.
\]
In particular the return is before $t_\eta$ for small $q$. The
continuous-inactivity hypothesis for transporting this band from
$t_h$ to weak learning is false. These quantifiers permit new
crossings later on times depending on $q$.
\end{theorem}

\section{A delay from physical coordinates}
Let $U,W$ be the balanced vectors, $\dd l=\dd\alpha/(2\pi)$,
$S=\int m\dd l$, and
$D_\perp=(\int(U_2^2+W_2^2)\dd l)^{1/2}$.
At $t_h$, HBD physical and clock give $S\le C$, $D_\perp\le Cq$.
For all later times P exactflow in vectors is
$U'=R_\theta^TW$, $W'=R_\theta U$,
$R_\theta=\E[(D-f)D^T{\bf1}_{U^TD>0}]$.
Balance and positive homogeneity give
$S'=2\E[f\cdot(D-f)]\le \E|D|^2/2$, hence $S\le C(1+s)$.
The exact training loss $\E|D-f|^2/2$ decreases along gradient
flow (its derivative is minus the integrated squared parameter
velocities); its initial value is uniformly bounded.
Cauchy--Schwarz then gives $\sup_\theta\|R_\theta\|\le C$.
For the off-diagonal entries one has the sharper bound
\[
 |R_{12}|+|R_{21}|\le C[q(1+S)+\sqrt S D_\perp].
\]
Indeed $|f_1(D)|\le S|D_1|+\sqrt S D_\perp|D_2|$,
$|f_2(D)|\le \sqrt S D_\perp|D|$, and $\E|D_1D_2|\le Cq$
on the two-cluster mixture. These bounds retain all donors for
every gate. Taking label $L^2$ norms in the second-component
equations gives, including by regularization at norm zero,
\[
 D^+D_\perp\le C(1+S)D_\perp+Cq\sqrt S(1+S).
\]
Integration yields $D_\perp\le C q\exp(C(1+s)^2)$ (enlarge $C$
to include the polynomial integral). Finally on $P_2$ itself,
$|f_2(D)|\le \sqrt S D_\perp |D_1|+D_\perp^2 |D_2|$.
Since $D_1$ has standard deviation $q$ on this cluster,
\begin{equation}\label{an04brwv1:delay}
 |\mathsf m_2|\le C(q\sqrt S D_\perp+D_\perp^2).
\end{equation}
Choose $k>0$ small enough for $C(1+k\sqrt{\log(1/q)})^2
\le \log(1/q)/2$ for small $q$ in the $D_\perp$ estimate.
Before $t_h$ the HBD physical bounds also imply
$|\mathsf m_2|\le C H^2(Eq+E^2)=O(q^2)$ uniformly, where
$E=e^{-\omega t}+q$. Indeed $H^2 e^{-2\omega t}
=S_0e^{(\lambda+2q^2)t}$ is increasing with endpoint $O(q^2)$.
This proves the delay, without giving an upper hitting time.

\section{The logistic resident reached from training}
HBD proves compact-time extraction, not a long-time $q$-uniform
angular approximation. We first extend its compact-time range.
For any fixed $T$, replace its nominal cap $S_0e^t=8$ by
$A_{\max}= 8 e^T$. In HBD physical the global envelope
$H^2=S_0e^{(1+2q^2)t}\le 2A_{\max}$ eventually; $\sqrt m\le
H\exp(H^2/2)$ and $\int H^2\,dt\le2A_{\max}$ still hold.
Integrating the same physical second-component inequality with
these finite constants bounds $\sup(|U_2|,|W_2|)$ by
$C_T H(e^{-\omega t}+q)$. HBD's small-cap left radial estimate
propagates across a fixed time by $|(\log m)'|\le(1+2q^2)(1+S)$,
and costs $o(q)$ in mass. A right interior compact keeps its
positive radial lower bound by the same inequality; the physical
second-component bound keeps its chart reached. Initial-label strips
of length $d$ cost $C_Td$ in scaled receiver velocities by precisely
the two Gaussian estimates of HBD. Hence its compact-kernel and
tagged-path argument applies on this larger interval with constants
depending on $T$. Diagonal extraction gives limiting paths for all
finite $s\ge0$, with the same HBD entrance and laws. None of these
constants are used as bounds at $s=t_\eta-t_h$.

In each such limit write $M=(1+e^{-s})^{-1}$,
$\dd\mu=M\dd\pi$, with fixed normalized right label weights
$\dd\pi=4\cos^2\alpha\,\dd l$. Conditional paths, if present,
use the same normalized marginal. With $P_x=\Phi(\rho x)$,
$K(z)=\phi(z)+z\Phi(z)$, $K_s=\int K(\rho x)\dd\mu$,
$Y=\int y\dd\mu$, and
$S_x=\int F(\rho x,\rho\tilde x)\dd\mu$ as in HBD,
\begin{equation}\label{an04brwv1:field}
 2x'=-(1-M)x+\rho\phi(\rho x)-\rho S_x+\lambda yP_x,\qquad
 2y'=-(1-M)y-Y+\rho K(\rho x).
\end{equation}
Here the weak limiting mass is zero at every finite $s$. The field
has nonnegative off-diagonal receiver derivatives. HBD gives a
single bottom Q4 path. Two different Q1 labels are ordered in both
coordinates because their remote incoming coefficients
$c^\omega\tan\alpha$ are strictly ordered in both coordinates
(HBD incoming); thereafter compare only within the common field.
The bottom path tends to $a_\rho>0$ at the remote end, so the
whole law has positive $Y$ sufficiently far back. Differentiation
gives $2Y'=(1-2M)Y+\rho K_s$; thus $Y>0$ at all finite times.
HBD's moment domination on each fixed interval justifies differentiation.

\paragraph{Uniform tail data at strong half mass.}
On Q1 let $f=\int_\alpha^{\pi/2}4\cos^2\beta\,\dd\beta/(2\pi)$.
It is the normalized upper mass. At the upper endpoint
$f\asymp(\pi/2-\alpha)^3$. HBD physical at half mass and
$m=2\cos^2\alpha$ imply
$x_+(0)+y_+(0)\le C(1+f^{-1/3})$, since each angular numerator
in the right-chart limit is bounded by
$2\sup(|U_2|,|W_2|)/q$ before division by $\sqrt m$.
The Q4 plateau is bounded using one interior representative.
Conversely HBD outer gives at half mass, for sufficiently high
Q1 labels,
$\min(x(0),y(0))\ge 4^{-\omega}(c^\omega\tan\alpha-J_*)$,
where $J_*$ depends only on the original HBD cap through $8$.
In particular the normalized mass with both coordinates initially
at least $R$ is at least $c_1(1+R)^{-3}$ for $R$ sufficiently
large, uniformly in the limiting law. These two bounds replace
the aligned unit-mass entry assumptions of IMP and FFP.

\paragraph{Two rates for different ranks.}
Set $L=\lambda/2$. For $s\ge0$, $M\ge1/2$. At a Q1 rank of upper
mass $f$, $S_x\ge MfK(\rho x)$ by order. Equation
\eqref{an04brwv1:field} gives the upper comparison
\[
 2D^+\binom{x_+}{y_+}\le
 \lambda\begin{pmatrix}-f/2&1\\1&0\end{pmatrix}
 \binom{x_+}{y_+}+\rho\phi(0)\binom11 .
\]
At negative coordinates the derivative of the positive part is
zero and the corresponding right side is nonnegative. At zero use
the positive-part derivative; on positive faces use $Y\ge0$ and
$K(r)\le r_++\phi(0)$. The same bound for bottom labels uses a
fixed upper mass, for instance $1/2$. Multiplying by the positive
left eigenvector $(w,1)$,
$w=(\sqrt{4+(f/2)^2}-f/2)/2$, and integrating gives
\[
 x_++y_+\le C(1+f^{-1/3})e^{Lws},\qquad
 Lw\le L-\lambda f/16.
\]
Use $\sqrt{4+d^2}\le2+d^2/4$ for $0\le d\le1$; $w$ stays bounded
away from zero. Since $\dd\pi=|\dd f|$ on Q1, integration of
$(1+f^{-1/3})e^{-\lambda fs/16}$ (substitute $u=fs$ when
$s\ge1$) proves the donor budget
\begin{equation}\label{an04brwv1:budget}
 \max(\rho K_s,Y)\le B_0 e^{Ls}(1+s)^{-2/3}.
\end{equation}
The constants are uniform in the possible extracted profiles;
only the half-mass tail data entered.

For each horizon $T$, let $g(s)=e^{Ls}\sqrt{M(0)/M(s)}$, so
$g\ge e^{Ls}/\sqrt2$ and $2g'/g=\lambda-(1-M)$.
Take $b_0> \sqrt2 B_0$ and a large fixed $d>0$. Put
\[
 R_T=d+3b_0[(1+T)^{1/3}-1],\qquad
 w_T(s)=g(s)\{R_T-3b_0[(1+s)^{1/3}-1]\}.
\]
Labels starting with both coordinates at least $R_T$ stay above
$w_T$ in both on $[0,T]$. On the $x=w_T,y\ge w_T$ face use
$\rho\phi(\rho x)+\lambda yP_x\ge\rho K(\rho w_T)$ and $S_x\le K_s$.
On $y=w_T,x\ge w_T$ use monotonicity of $K$. Each doubled velocity
strictly exceeds $2w_T'$ by \eqref{an04brwv1:budget},
$\rho K(\rho w_T)>\lambda w_T$ and $g\ge e^{Ls}/\sqrt2$.
This covers both faces and their corner. Therefore at every $s$ there
is a set, depending on $s$, of actual limiting radial mass
$\ge c_2/(1+s)$ with inputs $\ge d e^{Ls}/\sqrt2$.
Its first input moment is at least $c_3e^{Ls}/(1+s)$.

\paragraph{Return of the same labels.}
On the fixed band $B$, $f\ge f_0>0$; thus $x_+,y_+\le C_Be^{rs}$
for a single $r<L$. The protected set eventually lies above every
receiver in $B$. Its donors alone give
$S_x\ge\rho P_x c_3e^{Ls}/(1+s)$ because for $\tilde x\ge x$,
$F(\rho x,\rho\tilde x)=\phi(\rho x)+\rho\tilde xP_x$, and the
other overlap contributions are nonnegative. Fix $z=(a_\rho+\zeta)/2>0$.
Whenever $x\ge z$, \eqref{an04brwv1:field} now gives
\begin{equation}\label{an04brwv1:drift}
 2x'\le\rho\phi(0)+\lambda P_x[C_Be^{rs}-c_3 e^{Ls}/(1+s)]\le-1
\end{equation}
for $s\ge S_B$ independent of the extracted law, since
$P_x\ge\Phi(\rho z)>0$. At $S_B$ the heights are uniformly bounded
by the same ceiling. Integrate until first downward contact; the
strict sign at the gate prevents upward recrossing of $z$. There
is thus one finite $T_B$ after which all limiting $B$ paths satisfy
$x\le z$. With $d_B=(\zeta-z)/2$, strict separation from the actual
probe gate transfers on each fixed $[T_B,T]$ by tagged compactness
(otherwise extract a violating label and time). This completes the
proof of Theorem~\ref{an04brwv1:return}. For example at $T_B$ this
band even has a strictly positive horizontal probe output divided by
$q$ in the small-noise limit inferior: its limiting mass is positive,
$b_1\to1$, and $\sin(q(\zeta-x))/q$ is bounded below on this compact
band. Its zero probe-output assumption has a nonvanishing scaled error.

\paragraph{Trained survivor corollary.}\label{an04brwv1:survivors}
Weight each right-ancestry label by its mass at $t_h$.
The mass remaining probe-inactive at every time from $t_h$ to
$t_\eta$ tends to zero as $q\downarrow0$. If the hit never occurs,
take survival on the full infinite interval.
Indeed \eqref{an04brwv1:drift} and the return argument apply to
any fixed compact Q1 subinterval away from the vertical endpoint,
without initial inactivity being needed. Q4 interior compacts are
already active at $t_h$ by HBD. Discard right label strips of total
length $\delta$ about the axes; their half-mass weight is $O(\delta)$
by HBD's uniform per-label radial bound. For fixed $\delta$, all
remaining labels have been active by a fixed finite $T(\delta)$,
prior to $t_\eta$ for small $q$. Take the upper limit, then send
$\delta$ to zero. This concerns half-mass-time weights and continuous
survival, not current inactive mass or donor first moments at
$t_\eta$.

\section{What a learned-depth proof must pay}
\paragraph{Benchmark and clock.}
Use $(\rho,\zeta)$ in WIT's learned rectangle for this section.
Let $e_*=\min_{s\in J}\mathcal E_{\rm coh}(s)$, fix a minimizer
$s_m$, and write $A_m=\zeta-X_m>0$, $Z_m=X_m-Y_m>0$.
For the trained error put $e_q(t)=(E_\xi(t)-1/2)/q^2$.
For definiteness a measured comparison window can be centered on
$t_{1/2}$ as defined above, with test times $t\in t_{1/2}+J$
satisfying the required trained learning criteria. The chosen
response clock and the coherent logistic clock measure different
quantities; existence of such a reached learned window is not assumed.
The following identities work at any actual test time, hence also
at a learned minimizer if one exists. Beating $e_*$ on a learned
interval requires a time there with $e_q(t)<e_*$; a margin uniform
in $q$ gives strict asymptotic enhancement.

\paragraph{Exact accumulated test (trained flow).}
For a receiver-dependent matrix field $H$ define the scalar observable
\begin{align*}
 \mathcal P_t(H)=-q^{-2}\int\big[&
 2(b^THa)z_+(e_\xi^Tb)+(b_\perp^THa)z_+(e_\xi^Tb_\perp)\\
 &+(b^THa_\perp){\bf1}_{z>0}(a_\perp^T\xi)(e_\xi^Tb)\big]\dd\mu_t ,
\end{align*}
where $z=a^T\xi$, $e_\xi=\xi-f_t(\xi)$, $\dd\mu_t=m\dd l$,
and $a,b$ here are the unit input and output directions of each
current receiver. The Gaussian expectation in $H$ includes both
mixture weights. P motion gives $e_q'=\mathcal P_t(R)$ almost everywhere,
including radial motion and all current probe gates.
Fix a reference label in the HBD core interval and at each time
put $f_{\rm copy}(D)=\mu_t(B)b_C(a_C^TD)_+$. Define
\[
 H_B(a)=-\E[(f_B(D)-f_{\rm copy}(D))D^T{\bf1}_{a^TD>0}],
 \qquad H_{\rm rest}=R-H_B .
\]
This is the exact fixed-label donor-minus-copy attribution in the
trained field, with no inactivity restriction. It is linear at the
actual state and does not freeze the weak population. In particular
$\mathcal P_t(H_{\rm rest})$ includes motion of weak and strong
receivers, as does $\mathcal P_t(H_B)$.

To put the Step 1 cost on exactly the same scale, lift the coherent
state at $s_m$ to the finite-$q$ probe prediction
$F_c=b_{qY_m}\sin(qA_m)$ (unit active core, inactive weak atom).
Take $\varepsilon=\mu_{t_h}(B)/M_q(t_h)$, a mass fraction, and
$F_d=(1-\varepsilon)F_c$. Write
$e(F)=(|\xi-F|^2-1)/(2q^2)$, and define
\[
 C_q=e(F_d)-e(F_c)
 =q^{-2}\{\varepsilon(\xi-F_c)^TF_c+
                    \varepsilon^2|F_c|^2/2\}.
\]
This is an accounting reference, not an assertion that the trained
state ever equaled $F_d$. Uniformly for $0\le\varepsilon\le1$,
$C_q=\varepsilon A_m Z_m+\varepsilon^2 A_m^2/2+O(q^2)$
and $e(F_c)-e_*=O(q^2)$ by elementary trigonometric expansion.
Integration of P motion gives the following exact necessary and
sufficient strict test at any specified terminal time $t\ge t_h$:
\begin{align}
 -\int_{t_h}^{t}\mathcal P_u(H_B)\dd u
 >{}& C_q+[e_q(t_h)-e(F_d)]+[e(F_c)-e_*]\nonumber\\
    &+\int_{t_h}^t\mathcal P_u(H_{\rm rest})\dd u .
 \label{an04brwv1:payment}
\end{align}
Both sides use the actual residual and receiver weights. A comparison
to the exact lifted coherent minimum instead simply substitutes its
value for $e_*$. An angular-velocity integral over core receivers
alone omits the probe weights, radial reaction, other receivers and
the remainder on the right. It cannot verify this inequality.
Nor does this attribution give the difference between two evolved
training runs: such a counterfactual has its own residual and clock.

\paragraph{Evolved mass-matched comparison (trained flow).}
There is a stronger causal formulation of the unpaid term. Start a
second exact Gaussian flow at $t_0=t_h$ from the actual state with
only $B$ concentrated at the chosen active core coordinates, retaining
its radial mass. All other coordinates, including prospective weak
labels, agree at $t_0$. Evolve both systems. Superscripts $S,C$ denote
the factual split and concentrated runs; $C$ here is distinct from
the selected coherent orbit. Let $r^i(D)=D-f^i(D)$, $d=r^S-r^C$,
$\bar r=(r^S+r^C)/2$, and write $d_0=d(t_0)=f_{\rm copy}-f_B$.
At the probe $d_0(\xi)=f_{\rm copy}(\xi)$ by strict initial inactivity.
The physical removal cost in this experiment is
$C_0^{\rm act}=q^{-2}\bar r_{t_0}(\xi)^T d_0(\xi)$; it need
not equal the unit-core cost $C_q$.

For either state use its current kernel on vector $L^2(P)$,
\[
 K(z,D)=\int\!\left[(a^Tz)_+(a^TD)_+ I_2+
 (z^TD)\mathbf1_{a^Tz>0}\mathbf1_{a^TD>0}bb^T\right]\dd\mu .
\]
Write $K$ also for its integral operator and $k$ for its probe row
$K(\xi,\cdot)$. Differentiating the two vector factors of each ReLU
feature by P exactflow gives $(r^i)'=-K^i r^i$ on data, and
$(r^i(\xi))'=-k^i r^i$ almost everywhere. The kernels are bounded
operators on finite intervals by the mass and Gaussian second
moments; the selector chain rule also applies to the concentrated
atom. Let $\Delta K=K^S-K^C$, $\Delta k=k^S-k^C$.
Write $U(t,s)$ for the time-ordered propagator of $v'=-K^S v$
on the data space. For a finite terminal time $T$, define the map
from data residual to probe output
\[
 J_T(s)=\int_s^T k^S_t U(t,s)\dd t,\qquad
 B_T=J_T(t_0)d_0+
 \int_{t_0}^T[\Delta k_s-J_T(s)\Delta K_s]r^C_s\dd s .
\]
Indeed $d'=-K^S d-\Delta K r^C$. Substitute its
variation-of-constants solution into
$d'(\xi)=-k^S d-\Delta k r^C$ and interchange the two finite
integrals. This yields the exact identity
$d_T(\xi)=d_0(\xi)-B_T$.
Both direct probe-row changes and changes fed back through the
data residual are present. Positive semidefiniteness of $K^S$ on
the data gives no sign for the projected difference of these terms.

Define the scaled accumulated response
$\Pi_T=q^{-2}\bar r_T(\xi)^T B_T$. At $T$ the exact test for
beating the coherent minimum becomes
\begin{align}
 \Pi_T>{}&C_q+(C_0^{\rm act}-C_q)
 +q^{-2}[\bar r_T(\xi)-\bar r_{t_0}(\xi)]^Td_0(\xi)
 +e_q^C(T)-e_* .\label{an04brwv1:causal}
\end{align}
Here $e_q^C$ uses the independently evolved concentrated run.
The identity follows by polarizing squared norms at $T$.
For a test purely against that counterfactual, omit $e_q^C(T)-e_*$.
The first two terms on the right distinguish the actual half-mass
cost from Step 1's benchmark. The third accounts for the changing
probe residual weighting of the initially removed output. The
last compares the counterfactual's attained depth with the specified
benchmark. Even with the actual weak-learning time supplied,
the half-mass theorem gives no lower bound on $\Pi_T$ exceeding
this right side. In particular neither the sign of a local fitted
core velocity nor the disappearance of perpetual inactive ancestry
determines the sign of the memory term $J_T(s)\Delta K_s r_s^C$.

\paragraph{A second test using current output moments (trained flow).}
The following gives an endpoint target independent of the radial
and angular decomposition. Put $H=f_1(\xi)/q$, $V=f_2(\xi)/q^2$,
$s_q=\sin(q\zeta)/q$, $c_q=\cos(q\zeta)>0$, and
$h_1(a)=\E_{P_1}[D_1(a^TD)_+]/(1+q^2)$.
Let $I=\{a^T\xi\le0\}$ at the test time, $I^c$ its active complement,
and $\ell_G=\int_G b_2 h_1(a)\dd\mu$ for any current label set $G$.
The full regression is $r_{21}=\E_{P_1}[D_1(f_2-D_2)]/(1+q^2)
=\int b_2h_1(a)\dd\mu$. For fixed $k>0$ define
\[
 G_k=q^{-2}\int_{I^c} b_2[a^T\xi-qk h_1(a)]\dd\mu .
\]
P transferprojection, with all labels, gives
$V=k(r_{21}-\ell_I)/q+G_k$. The exact squared probe loss gives
$e_q=H^2/2-s_qH+c_q V+q^2V^2/2$. Thus the same depth claim is
equivalent at the test time to
\begin{equation}\label{an04brwv1:endpoint}
 \frac{c_q k}{q}\ell_{I\cap B}>
 \frac{H^2}{2}-s_qH+c_qG_k+
 \frac{c_q k}{q}(r_{21}-\ell_{I\setminus B})+
 \frac{q^2V^2}{2}-e_* .
\end{equation}
This test is valid even if $I\cap B$ is empty; earlier influence
can then remain in $G_k$, the regression, and other donors. It neither
turns an active band into an inactive regression term nor drops the
quadratic output penalty. The falsified guard is
$x_\alpha\ge\zeta$ at all intervening times for $\alpha\in B$,
with the exact finite-noise guard $a_\alpha^T\xi\le0$.
The trained estimates through the pre-learning interval do not pay
\eqref{an04brwv1:payment} or \eqref{an04brwv1:causal}, or bound the right side of
\eqref{an04brwv1:endpoint} below the left. Neither inequality
is shown false at a learned minimum.

\paragraph{Source ledger.}
HBD denotes \texttt{an04hbdv1:} main, clock, physical, incoming,
outer and field, from the local trained half-mass increment.
FFP return and IMP mass and budget concern unit strong mass and
their prescribed aligned entrance. Here the changing logistic
mass, upper-moment ceiling and lower barrier were rederived for
the HBD entrance. Step 1 denotes \texttt{an04mpbv1:depth}, for
a prescribed short split at unit strong mass with the weak atom
retained. P's \texttt{pa:exactflow}, \texttt{pa:learning},
\texttt{pa:motion}, and \texttt{pa:transferprojection} are exact
Gaussian trained identities under finite moments and characteristic
evolution with the probe chain rule a.e. WIT uniformtheorem supplies
only the coherent learned-window benchmark. No finite-$q$ empirical
trajectory or proposed weak-seed law enters a bound.
\end{document}
